\documentclass{ltjsarticle}

\setlength\parindent{0pt}
\setlength{\columnseprule}{0.4pt}

\newcommand{\sk}{\vskip\baselineskip}

\usepackage[margin=15mm]{geometry}
\usepackage{amsmath}
\usepackage{amssymb}
\usepackage{amsfonts}
\usepackage{mathrsfs}
\usepackage{bm}
\usepackage{here}
\usepackage{bbm}
\usepackage{graphicx}
\usepackage[unicode]{hyperref}
\usepackage{mathtools}
\usepackage{url}
\usepackage{ascmac}
\usepackage{physics}
\usepackage[version=4]{mhchem}
\usepackage{listings}

\lstset{
language=Python, 
basicstyle={\small},
identifierstyle={\small},
commentstyle={\small\itshape},
keywordstyle={\small\bfseries},
ndkeywordstyle={\small},
stringstyle={\small\ttfamily},
frame={tblr},
breaklines=true,
columns=[l]{fullflexible},
xrightmargin=0zw,
xleftmargin=3zw,
numberstyle={\scriptsize},
stepnumber=1,
numbersep=1zw,
lineskip=-0.5ex
}

\mathtoolsset{showonlyrefs=true}

\title{ランダム行列 常識集}
\author{東京大学 理学系研究科物理学専攻 修士1年 平松大門}
\begin{document}
\maketitle
\section{Getting Started}
\begin{itemize}
    \item Gaussian Ensemblesの代表例3つは？
    \begin{table}[h]
\centering
\begin{tabular}{|c|c|c|c|}
\hline
略称 & 正式名称&行列の性質&Dyson Index $\beta$\\
\hline
GOE & Gaussian Orthogonal Ensemble & 実対称行列\quad $H=H^T$ &1
\\
\hline
GUE & Gaussian Unitary Ensemble&エルミート行列 \quad$H=H^\dagger$ &2
\\
\hline
GSE &Gaussian Symplectic Ensemble& 四元数自己双対エルミート行列 &4
\\
\hline
\end{tabular}
\end{table}

Dyson Indexは行列の成分が持つ実自由度の数を表す。

\sk
\item 重積分の変数変換は？
\[
\int\prod_i dx_i \,f(x_1, \cdots ,x_n) = \int \prod_i dy_i\, \left|\frac{\partial (x_1, \cdots , x_n)}{\partial (y_1, \cdots ,y_n)}\right|f\left(x_1(\bm{y}), \cdots, x_n(\bm{y})\right)
\]

\end{itemize}

\newpage

\section{Value the Eigenvalue}
\begin{itemize}
    \item Wigner予想とは？
    
    ランダム行列の固有値間隔の分布についての予想。$2\times 2$のGOEの固有値間隔$s$の分布$P(s)$を期待値が1になるようリスケールした分布$\tilde{P}(s)$は
    \[
    \tilde{P}(s)=\frac{\pi}{2}s\exp(-\frac{\pi}{4}s^2)
    \]
    となるが、$N\times N$でも似た分布になると予想した。
    固有値間隔が重なる確率が0、つまり固有値同士は反発することを意味する。
\end{itemize}


\end{document}